% Positive effective geometry for lattice-polarised K3 x E families.
% Primary sources: Baily-Borel 1966; Dolgachev 1996; Nikulin 1979, 1980;
% Borcherds 1998; Gritsenko-Nikulin 1995, 1998; Gritsenko 1999;
% Clery-Gritsenko 2008; Gritsenko-Hulek-Sankaran 2007; Mukai 1987, 1988;
% Maulik-Pandharipande-Thomas 2010; Kiem-Li 2013; Oberdieck 2018.

\providecommand{\Z}{\mathbb{Z}}
\providecommand{\C}{\mathbb{C}}
\providecommand{\bP}{\mathbb{P}}
\providecommand{\gfrak}{\mathfrak{g}}
\providecommand{\gL}{\mathfrak{g}_{L}}
\providecommand{\Pic}{\mathrm{Pic}}
\providecommand{\Om}{\Omega}
\providecommand{\wt}{\mathrm{wt}}
\providecommand{\rk}{\mathrm{rk}}
\providecommand{\Oneg}{\mathrm{O}^{+}}
\providecommand{\cM}{\mathcal{M}}
\providecommand{\cX}{\mathcal{X}}
\providecommand{\cO}{\mathcal{O}}
\providecommand{\LamK}{\Lambda_{K3}}
\providecommand{\Mred}{\mathcal{M}^{\mathrm{red}}}
\providecommand{\CoHA}{\mathrm{CoHA}}
\providecommand{\Rep}{\mathrm{Rep}}
\providecommand{\End}{\mathrm{End}}
\providecommand{\Perf}{\mathrm{Perf}}
\providecommand{\kBKM}{\kappa_{\mathrm{BKM}}}
\providecommand{\kch}{\kappa_{\mathrm{ch}}}
\providecommand{\kcat}{\kappa_{\mathrm{cat}}}
\providecommand{\kfib}{\kappa_{\mathrm{fiber}}}

\section*{The lattice-polarised family
  $\gL/\cM_{L}$: positive effective geometry}
\label{sec:gL-positive-geom-family}

\paragraph{Thesis.}
Let $L$ be an even non-degenerate lattice of signature $(1,t)$.
Assume that $L$ is Nikulin-Dolgachev admissible.  Then the open moduli
stack of strictly $L$-polarised K3 surfaces carries, after passage to a
neat finite level cover when a universal surface is required, a family
\[
  \cX_L\times E\longrightarrow \cM_L^\circ
\]
of Calabi-Yau threefolds.  Primitive reduced Donaldson-Thomas theory on
the K3 factor supplies the coefficient package that enters the
corresponding Borcherds product.  The BKM invariant is the automorphic
weight
\[
  \kBKM(\gL)=\wt(\Phi_L)=\frac{c_L(0)}{2},
\]
where $c_L(0)$ is the zeroth Fourier coefficient of the chosen
Borcherds input in the stated normalisation.  This is denominator data.
The Mukai Yangian branch is a separate K3 surface construction, and the
full Hall-Drinfeld double requires a positive half, a negative half, a
Cartan factor, a non-degenerate Hall pairing, completions, and descent.
For $d=3$ the native output of $\Phi_3$ is $E_1$-chiral; $E_2$-braiding
belongs to the Drinfeld centre or to the completed double.

\subsection*{Nikulin-Dolgachev admissibility}

The K3 lattice is
\[
  \LamK=2E_8(-1)\oplus 3U .
\]
An even lattice $L$ of signature $(1,t)$ is admissible for this
construction when
\[
  L\hookrightarrow \LamK
  \quad\text{is primitive},
  \qquad
  L^\perp\supset U .
\]
The primitive embedding is governed by Nikulin's discriminant-form
criterion (Nikulin 1980, Theorem 1.12.2).  Writing
$A_L=L^\vee/L$, the condition is a finite arithmetic condition on the
length, parity, and local discriminant form of $A_L$.  The second
condition is Dolgachev's mirror condition:
\[
  L^\perp=U\oplus\check L .
\]
Since $\rk L=t+1$, the orthogonal complement has rank $21-t$, and
\[
  \rk\check L=19-t,\qquad
  \mathrm{sign}(\check L)=(1,18-t).
\]
Primitive embeddability and the mirror split are different hypotheses.
All lattice towers used below are required to remain inside this
admissible locus.

\subsection*{Period domain and universal family}

Put $T=L^\perp$, of signature $(2,19-t)$.  The period domain is
\[
  \Om_L=
  \bigl\{[\omega]\in\bP(T\otimes\C):
    \langle\omega,\omega\rangle=0,
    \langle\omega,\overline{\omega}\rangle>0\bigr\}^{+}.
\]
It is a Type IV Hermitian symmetric domain of dimension $19-t$
(Dolgachev 1996, Theorem 3.1).  The arithmetic group
$\Gamma_L=\Oneg(T)$ acts properly discontinuously, and
\[
  \cM_L=\Gamma_L\backslash\Om_L
\]
is the coarse moduli space of $L$-polarised K3 surfaces.  The moduli
stack, or a neat finite level cover of the quotient, carries a
tautological K3 surface
\[
  \pi_L:\cX_L\longrightarrow \cM_L .
\]
The strict open locus $\cM_L^\circ$ is the complement of the
Noether-Lefschetz hyperplane arrangement in which a class of $T$ becomes
of type $(1,1)$.  On this locus
\[
  \Pic(X)=\iota_L(L)
\]
for a very general point, and the embedding $\iota_L:L\hookrightarrow
\Pic(X)$ remains primitive throughout the family.  With $E$ a fixed
smooth elliptic curve, the product family
\[
  \cX_L\times E\longrightarrow \cM_L^\circ
\]
has fibres $X\times E$ with trivial canonical bundle and
$\kcat(X\times E)=0$.

\subsection*{Reduced DT coefficients and the Borcherds divisor}

Fix a primitive effective class $\gamma\in L$ with
$\gamma^2=2h-2$.  For $Y=X\times E$, the primitive curve-counting
problem is the reduced stable-pair or ideal-sheaf theory in class
$\gamma$ and Euler parameter $n$, with the translation action of $E$
divided out.  The raw obstruction theory has the K3 symplectic
cosection
\[
  \sigma_L^*:\mathrm{Ob}\longrightarrow\cO
\]
from the holomorphic two-form on the K3 factor.  Kiem-Li cosection
localisation and the reduced K3 obstruction theory remove this trivial
direction; after the elliptic quotient the fibrewise reduced invariant is
zero-dimensional:
\[
  N_L^{\mathrm{red}}(\gamma,n)
  =
  \int_{
    [\Mred_{\mathrm{DT}}(Y,\gamma,n)/E]^{\mathrm{vir,red}}
  }1 .
\]
This is the numerical coefficient used in the reduced DT/PT partition
series of Maulik-Pandharipande-Thomas and Oberdieck.

The divisor $D_{\Phi_L}$ is the Borcherds divisor of the automorphic
product $\Phi_L$ on the Type IV domain, pulled back by the period map.
The reduced DT series and the Borcherds product agree only after fixing
the primitive chamber, orientation, elliptic quotient convention, and
Baily-Borel cusp normalisation.  In that normalisation the scalar
comparison is a theorem in the $K3\times E$ primitive branch; the
construction of a compact Hall-Drinfeld double is an additional algebraic
problem.

\subsection*{Borcherds packages and constants}

For a Borcherds product $\Phi$ with input coefficient $c_\Phi(0)$,
Borcherds' weight theorem gives
\[
  \wt(\Phi)=\frac{c_\Phi(0)}{2}.
\]
The BKM invariant records this weight:
\[
  \kBKM(\Phi)=\wt(\Phi)=\frac{c_\Phi(0)}{2}.
\]
The formula is rowwise.  The source package, divisor geometry, and
weight convention must be named before values are compared.

\paragraph{Primitive CHL ladder.}
For the CHL diagonal quotient package on $K3\times E$, with
$N\in\{1,2,3,4,6\}$ and $\varphi(N)\mid 2$, the Borcherds input is the
normalised twined weight-zero index-one Jacobi form used by
Gritsenko-Nikulin and Gritsenko.  The resulting paramodular forms
$\Phi_N$ have Borcherds divisors prescribed by their principal parts;
at $N=1$ this is Gritsenko's $\Delta_5$ on the diagonal Humbert
divisor.  The convention is
\[
\begin{array}{c|ccccc}
N & 1 & 2 & 3 & 4 & 6\\ \hline
c_N(0) & 10 & 8 & 6 & 4 & 2\\
\kBKM(\Phi_N)=c_N(0)/2 & 5 & 4 & 3 & 2 & 1 .
\end{array}
\]
The physics Igusa convention squares the form, so the physical weight is
$c_N(0)$ while the BKM denominator weight is $c_N(0)/2$.

\paragraph{Nikulin orders outside the CHL quotient.}
The symplectic K3 orders $N\in\{5,7,8\}$ use the Nikulin twined K3
input and extended automorphy groups.  They are separate from the CHL
quotient geometry:
\[
\begin{array}{c|ccc}
N & 5 & 7 & 8\\ \hline
c_N(0) & 4 & 2 & 2\\
\kBKM(\Phi_N)=c_N(0)/2 & 2 & 1 & 1 .
\end{array}
\]
Together with the CHL rows this gives the Nikulin-order sequence
\[
  c_N(0)=(10,8,6,4,4,2,2,2),\qquad
  \kBKM=(5,4,3,2,2,1,1,1)
\]
for $N=(1,2,3,4,5,6,7,8)$.  The entries $5,7,8$ have Borcherds weights;
their smooth $K3\times E$ CHL source package would require an
origin-fixing elliptic automorphism of order $5$, $7$, or $8$, which the
elliptic automorphism classification excludes.

\paragraph{Gritsenko-Clery diagonal-divisor catalogue.}
Gritsenko-Clery classify dd-modular forms whose divisor is the diagonal
$\cH_1$ with vanishing order $1$ on a paramodular quotient.  Their row
index is a triple $(t,N;k)$, where $k$ is the modular weight:
\[
\begin{array}{c|c|c|c|c}
R & (t,N;k) & F_R & c_R(0) & \kBKM(F_R)\\ \hline
1 & (1,1;5) & \Delta_5 & 10 & 5\\
2 & (2,1;2) & \Delta_2 & 4 & 2\\
3 & (1,2;3) & \nabla_3 & 6 & 3\\
4 & (3,1;1) & \Delta_1 & 2 & 1\\
5 & (1,3;2) & \nabla_2 & 4 & 2\\
6 & (4,1;\tfrac12) & \Delta_{1/2} & 1 & \tfrac12\\
7 & (1,4;\tfrac32) & \nabla_{3/2} & 3 & \tfrac32\\
8 & (2,2;1) & Q_1 & 2 & 1 .
\end{array}
\]
Thus
\[
  \wt(F_R)=(5,2,3,1,2,\tfrac12,\tfrac32,1),
  \qquad
  c_R(0)=(10,4,6,2,4,1,3,2).
\]
Rows $6$ and $7$ use multiplier systems of orders $8$ and $4$.  This is
a diagonal-divisor catalogue.  Its only common row with the primitive
CHL ladder is $\Delta_5$.

\subsection*{Quasi-pullback base change}

For a primitive admissible embedding $L\hookrightarrow L'$, the period
domain of $L'$-polarised surfaces is the subdomain of $\Om_L$ cut out by
the additional algebraic classes.  Borcherds products on the two domains
can be compared only by quasi-pullback.  The required data are:
\[
  \text{root set perpendicular to }\Om_{L'},\quad
  \text{Weyl chamber},\quad
  \text{reflection character},\quad
  \text{cusp normalisation}.
\]
After these choices the pullback is a Borcherds product on the smaller
domain multiplied by explicit root factors, and the weight changes by
one half of the total multiplicity of the removed perpendicular roots.
The simplified identity
\[
  \kBKM(\gfrak_{L'})-\kBKM(\gL)
  =
  \frac{c_{L'}(0)-c_L(0)}{2}
\]
is valid only for a certified descent between two named Borcherds
packages.  The chosen quasi-pullback supplies the identity; arbitrary
lattice towers fall outside a formal functorial weight law.

\subsection*{Positive half, centre, and chiral output}

The local toric normalisation is the Schiffmann-Vasserot identity
\[
  \CoHA(\C^3)=Y^+(\widehat{\mathfrak{gl}}_1),
\]
where the right side is the positive half of the affine Yangian.  The
full affine Yangian and the $\mathcal W_{1+\infty}[\lambda]$ vertex
algebra appear after adjoining the Cartan and negative halves, completing
the Hopf pairing, and evaluating:
\[
  Y^+(\widehat{\mathfrak{gl}}_1)
  \hookrightarrow
  D\bigl(Y^+(\widehat{\mathfrak{gl}}_1)\bigr)
  =
  Y^+\otimes Y^0\otimes Y^-
  \xrightarrow{\mathrm{ev}_{\lambda}}
  \End\bigl(\mathcal W_{1+\infty}[\lambda]\text{-vac}\bigr).
\]
For the compact family $\cX_L\times E$, positive effective geometry
supplies Hall-positive and Borcherds-positive data.  The centre-side
object is
\[
  \mathcal Z\bigl(\Rep^{E_1}(A_Y)\bigr),
  \qquad
  A_Y=\Phi_3^{(\Sigma_2,C)}(\Perf(Y)),
\]
or the corresponding completed Hall-Drinfeld double after pairing and
completion.  Hence the $d=3$ $\Phi$ output is $E_1$-chiral; $E_2$
structure belongs to the centre, double, or module category.  The BKM
denominator branch and the Mukai Yangian branch share K3 input but use
different objects, different lattices, and different invariants.

\subsection*{Invariant separation on \(K3\times E\)}

For $Y=X_L\times E$ the relevant invariants are:
\[
\begin{array}{c|c|l}
\text{invariant} & \text{value} & \text{source}\\ \hline
\kcat(Y) & 0 &
  \chi(\cO_{X_L})\chi(\cO_E)=2\cdot0\\
\kch^{\mathrm{Hodge}}(Y) & 0 &
  \sum_q(-1)^q h^{0,q}(Y)=1-1+1-1\\
\kch^{\mathrm{Heis}}(Y) & 3 &
  \text{Heisenberg-Mukai chiral specialisation}\\
\kBKM(\mathfrak g_{\Delta_5}) & 5 &
  c_1(0)/2=10/2\\
\kfib(Y) & 24 &
  \rk\,\widetilde H(K3,\Z)=1+22+1 .
\end{array}
\]
The value $2=\chi(\cO_{K3})$ is a fibre holomorphic Euler
characteristic.  The value $24$ is the K3 Mukai-lattice rank, equivalently
the topological Euler characteristic of the K3 fibre.  No additive law
relates these numbers:
\[
  \kBKM(\Phi_N)=\frac{c_N(0)}{2}
\]
is the complete BKM formula in its denominator scope.

\paragraph{Primary references.}
Baily-Borel 1966 \emph{Ann. Math.} 84;
Borcherds 1998 \emph{Invent. Math.} 132;
Clery-Gritsenko 2008 \texttt{arXiv:0802.3671};
Dolgachev 1996 \emph{J. Math. Sci.} 81;
Gritsenko 1999 \emph{St. Petersburg Math. J.} 11;
Gritsenko-Hulek-Sankaran 2007 \texttt{arXiv:math/0605380};
Gritsenko-Nikulin 1995 \texttt{arXiv:alg-geom/9504006};
Gritsenko-Nikulin 1998 \emph{Int. J. Math.} 9;
Kiem-Li 2013 \emph{J. Amer. Math. Soc.} 26;
Maulik-Pandharipande-Thomas 2010 \texttt{arXiv:1001.2719};
Mukai 1987 \emph{Nagoya Math. J.} 106;
Mukai 1988 \emph{Invent. Math.} 94;
Nikulin 1979 \emph{Izv. Akad. Nauk SSSR} 43;
Nikulin 1980 \emph{Tr. Mosk. Mat. Obs.} 38;
Oberdieck 2018 \texttt{arXiv:1605.05238}.
